% ~600 words. Submission, execution, aggregation. Comparator table.
\section{Protocols}
\label{sec:protocols}

\begin{figure*}[t!]
  \centering
  \resizebox{\textwidth}{!}{%
  \begin{tikzpicture}[
    font=\sffamily\small, >=Stealth,
    node distance=7mm and 9mm,
    card/.style 2 args={rounded corners=3pt, draw=#1!55!black, fill=#1!10,
      line width=0.6pt, text width=#2, inner sep=4pt, minimum height=15mm, align=left},
    badge/.style={circle, fill=#1!70!black, text=white,
      font=\sffamily\bfseries\footnotesize, inner sep=1.3pt, minimum size=5mm},
    zone/.style={rounded corners=6pt, draw=#1!55!black, dashed, line width=0.9pt, fill=#1!5},
    flow/.style={-{Stealth[length=3mm]}, line width=1.2pt, draw=#1!70!black},
    lbl/.style={fill=white, inner sep=1.5pt, font=\sffamily\footnotesize},
  ]
  % ---- top band: central develop + distribute (blue) ----
  \node[card={blue}{37mm}] (n1)
    {\textbf{FLAIR Library}\\ \emph{central $\cdot$ public}\\[1pt]
     \texttt{flair-baseline train}\\ cohorts $\cdot$ 75/25 split\\
     count feats $\cdot$ XGBoost\\ $\rightarrow$ \texttt{model.json}, \texttt{vocab.json}};
  \node[card={blue}{37mm}, right=of n1] (n2)
    {\textbf{Build / submit model}\\ baseline \emph{or} community PR\\[1pt]
     Benchmark Committee\\ review (clinical $\cdot$ security)\\ package for release};
  \node[card={blue}{37mm}, right=of n2] (n3)
    {\textbf{Distribute}\\ FLAIR Portal ships the\\ approved model $+$ code\\
     to all 12+ CLIF sites};
  \node[badge={blue}] at (n1.north west) {1};
  \node[badge={blue}] at (n2.north west) {2};
  \node[badge={blue}] at (n3.north west) {3};

  % ---- middle band: site silos (orange) ----
  \node[card={orange}{37mm}, below=13mm of n1] (s1)
    {\textbf{Site 1}\\ local CLIF tables\\ \texttt{flair-baseline infer}\\
     $\rightarrow$ report bundle ($n{<}10$)};
  \node[card={orange}{37mm}, right=of s1] (s2)
    {\textbf{Site 2}\\ local CLIF tables\\ \texttt{flair-baseline infer}\\
     $\rightarrow$ report bundle ($n{<}10$)};
  \node[card={orange}{37mm}, right=of s2] (sN)
    {\textbf{Site N}\\ local CLIF tables\\ \texttt{flair-baseline infer}\\
     $\rightarrow$ report bundle ($n{<}10$)};
  \node[font=\sffamily\bfseries, text=orange!45!black] (sdots) at ($(s2.east)!0.5!(sN.west)$) {$\cdots$};
  \begin{scope}[on background layer]
    \node[zone={orange}, fit=(s1)(sN), inner sep=7mm] (silo) {};
  \end{scope}
  \node[anchor=north west, font=\sffamily\bfseries\footnotesize, text=orange!50!black]
    at ([yshift=6.4mm]silo.north west) {Patient data never leaves the silo $\cdot$ privacy boundary $\cdot$ 12+ sites};

  % ---- bottom band: aggregate + publish ----
  \node[card={blue}{40mm}, below=14mm of silo.south west, anchor=north west] (a4)
    {\textbf{Aggregate}\\ \emph{central FLAIR}\\ pool the per-site\\ report bundles};
  \node[card={violet}{56mm}, right=of a4] (a5)
    {\textbf{Publish $\cdot$ leaderboard}\\ per-site \textbf{forest plot} (the headline)
     $+$ pooled scalar\\ HIPAA-safe $\cdot$ no PHI shared};
  \node[badge={blue}] at (a4.north west) {4};
  \node[badge={violet}] at (a5.north west) {5};

  % ---- arrows ----
  \draw[flow={blue}] (n1) -- (n2);
  \draw[flow={blue}] (n2) -- (n3);
  \draw[flow={blue}] (n3.south) to[out=-90,in=0]
    node[lbl, pos=0.62] {model $+$ code ships $\blacktriangleright$} (silo.east);
  \draw[flow={orange}] (silo.west) to[out=180,in=90]
    node[lbl, left, pos=0.32] {$\blacktriangleleft$ aggregate JSON $\cdot$ no PHI} (a4.north);
  \draw[flow={violet}] (a4) -- (a5);
  \end{tikzpicture}}
  \caption{\textbf{The FLAIR model-to-data loop.} Patient-level data never
  leaves any site; only a model crosses in and an aggregate report crosses
  out. \emph{Develop (central, public):} a model is built against MIMIC-CLIF
  with \texttt{flair-baseline train} (cohorts, a 75/25 split on
  \texttt{hospitalization\_join\_id}, count features, an XGBoost head) or
  submitted by the community as a pull request; the Benchmark Committee
  reviews it for clinical and security soundness. \emph{Ship:} the FLAIR
  Portal distributes the approved model folder and code to the 12+ CLIF
  sites. \emph{Run locally:} each Site Lead / Data Steward applies it with
  \texttt{flair-baseline infer} on the site's own CLIF tables, which writes
  three site-prefixed folders, and returns only the non-PHI \textbf{report
  bundle}, a \emph{mode-keyed} set of JSONs (\texttt{discrimination},
  \texttt{calibration}, \texttt{dca}, \texttt{fairness}, and, for landmark
  tasks, \texttt{leadtime}) plus \texttt{table1.json} and a suppressed code
  registry, with cell suppression at $n<10$ --- metrics only, no patient
  rows. \emph{Aggregate \& publish:} the Portal pools the per-site reports
  into the public leaderboard, where the per-site \textbf{forest plot}, not
  the pooled scalar, is the headline signal.}
  \label{fig:workflow}
\end{figure*}

The federated protocol is what makes FLAIR's privacy properties
operational. We first name the actors that operate the ecosystem, then
describe the protocol in three steps, developer-side submission,
site-side execution, and central aggregation, and position it
against the closest precedents.

\subsection*{The FLAIR ecosystem: roles and the Portal}

FLAIR is operated as an ecosystem of four roles, mirroring the
model-to-data benchmarking pattern MedPerf established for medical
imaging \citep{karargyris2023medperf} and bringing it to critical-care
longitudinal EHR. The \emph{FLAIR Portal} is the orchestration layer: it
registers tasks, models, and sites, distributes the containerised code
and Dockerfiles, collects the returned aggregate JSONs, and renders the
leaderboard. A \emph{Benchmark Committee}, clinicians, methodologists,
regulators, and patient representatives, reviews each new task, which
enters as a single-file Python module via GitHub pull request, for
reproducibility against MIMIC-CLIF and clinical face validity. The same
Committee is the privacy and safety guarantee on the model side: before a
submission is released to sites it security-reviews the submitted code and
model and validates the task code against synthetic (fake) data, so only
audited artifacts are ever dispatched, whether shipped as raw code or
wrapped in a container. A \emph{Model Owner} submits a frozen model, as
code or wrapped in a container. At each institution a \emph{Site Lead / Data
Steward} pulls the approved submission from the Portal, runs it in
CLIF-WorkBench against local CLIF tables, and returns only the
standardized non-PHI report bundle. The lifecycle is therefore
register~$\rightarrow$~approve~$\rightarrow$~execute~$\rightarrow$~release:
no patient-level data leaves any site, and every site runs byte-identical
Portal-distributed code.

\subsection*{Submission}

A clinical decision support developer using FLAIR works entirely
against MIMIC-CLIF during development. The library exposes CLIF tables
through \texttt{clifpy}, the canonical CLIF Python interface. The
benchmark CLI exposes \texttt{flair load-task} (build a task-specific
cohort parquet) and \texttt{flair build-report} (score predictions into
the report bundle); the reference baseline wraps these as
\texttt{flair-baseline train} and \texttt{flair-baseline infer}. Each
run writes three site-prefixed folders, separating PHI (cohort, events,
predictions; stays local) from the non-PHI upload bundle and the model
folder that ships to other sites. Submission to the consortium consists
of code and a frozen model artifact, optionally packaged as a container
image (Docker or Singularity) that wraps both. The model-to-data paradigm
itself is not new, it was introduced by the DREAM Challenges
\citep{guinney2018dream} and demonstrated for EHR-based mortality
prediction at scale in the EHR DREAM Challenge
\citep{bergquist2024ehrdream}. FLAIR's contribution is to move it from
a one-time competition format into a standing benchmark and submission
service. Before any model is dispatched to a site, the Portal runs its
container against MIMIC-CLIF as a compatibility check, so only models
that execute cleanly on the schema-locked public data are approved for
federated execution.

\subsection*{Execution}

Each consortium site runs the committee-reviewed, security-approved
submission, raw code or a container, inside the CLIF-WorkBench
environment. Because every submission has already passed central security
review of its code and model and a synthetic-data validation pass before
release, sites execute only audited artifacts. The submitted code is
permitted to read the site's CLIF tables, write a predictions parquet (cohort plus
\texttt{y\_true}, \texttt{y\_pred}, \texttt{y\_prob}), and invoke the
FLAIR report builder, which produces a mode-keyed report bundle
(\texttt{discrimination}, \texttt{calibration}, \texttt{dca},
\texttt{fairness}, and \texttt{leadtime} for landmark tasks; peak tasks
omit calibration and decision curves). Every JSON carries the FLAIR
library version, the generation timestamp, and the task audit hash. Sites retain final approval over which models
execute; the simplicity of the per-task module, one Python file
expressing cohort, label, and exclusions in a few hundred lines, 
makes that pre-execution audit tractable.

\subsection*{Aggregation, leaderboard, and lifecycle}

The central aggregator merges the per-site aggregate JSONs for a given
task version and publishes per-site and pooled metrics. Per-site variance is the primary signal of
generalization; pooled metrics are reported but de-emphasized. The
leaderboard records both retrospective benchmark performance and, for
tasks where sites have opted in, prospective silent-trial evaluation
under emerging guidance for clinical AI preregistration
\citep{goodman2025preregistration,joshi2025silenttrials}. Recent work
on operational templates for monitoring AI in routine clinical practice
\citep{pragmaticplaybook2024} informs the prospective arm.

\paragraph{Comparison with adjacent federated infrastructures.}
The closest precedent for federated multi-site evaluation in healthcare
is the OHDSI network, whose LEGEND framework standardized network
studies on the OMOP common data model \citep{reps2020ohdsistroke};
LEGEND federates SQL queries and coefficient sharing for logistic
models. MedPerf \citep{karargyris2023medperf} demonstrates federated
benchmarking for medical imaging models packaged as containers; the
FeTS Challenge \citep{pati2022fets} established federated training,
rather than validation, as a viable mode for imaging. OpenSAFELY
\citep{opensafely2021naturemedicine} delivers a code-travels-not-data
analytic platform for national-scale primary-care data. FLAIR
contributes the missing combination for critical-care EHR-based CDS:
modern longitudinal model classes, a clinically scoped data
schema (CLIF) shared across sites, and a community-extensible task
catalog under open governance. The five-attribute comparison is in
Table~\ref{tab:comparators}.

\begin{table}[t]
\centering
\footnotesize
\setlength{\tabcolsep}{3pt}
\caption{FLAIR alongside adjacent federated infrastructures.}
\label{tab:comparators}
\begin{tabularx}{\linewidth}{lXXXXX}
\toprule
 & OHDSI/LEGEND & MedPerf & FeTS & OpenSAFELY & \textbf{FLAIR} \\
\midrule
Data model       & OMOP         & task-specific & imaging (brain MRI) & TPP / EMIS       & CLIF \\
Federation unit  & SQL / coefficients & container & training gradients & analytic code & container / model code \\
Modern DL CDS    & no           & yes (imaging) & yes (training only)   & limited        & yes (longitudinal EHR) \\
Hosted-model portal & no & yes & yes & no & yes \\
PR-extensible tasks & no        & partial       & no                    & n/a             & yes \\
Sites (current)  & many         & \est{$\sim$10s} & 71                  & English NHS    & 12+ \\
\bottomrule
\end{tabularx}
\end{table}
